% Pista ana-26j2-q4a · Anàlisi
% Procedència: PAU Catalunya 2026 · Convocatòria de juny · Sèrie 5 · Exercici 4 — Opció A
\documentclass[11pt,a4paper]{article}
\usepackage{pista}
\pistainfo{Catalunya 2026}{Convocatòria de juny}{Sèrie 5}

\begin{document}

\section*{\textcolor{blauPista}{Exercici 4 --- Opció A}}

\noindent D'una funció $f(x)$, sabem que passa pel punt $(-1, -1)$ i que la gràfica de la seva funció derivada és la següent:

\noindent\textit{(A la figura: la gràfica de $f'(x)$, que val 0 a $x = -2$, $x = 0$ i $x = 1$.)}

\subsection*{4.1. \normalfont Calculeu l'equació de la recta tangent a la gràfica de la funció $f(x)$ en el punt $(-1, -1)$. \hspace*{\fill}\textnormal{[1~punt]}}

\pista{El pendent de la recta tangent que busques és $f'(-1)$, i el punt és $(-1, -1)$. Per tant:
\[
y - f(-1) = f'(-1)\,(x + 1)
\]
Fixa't que la gràfica de l'enunciat és la d'$f'$, no la d'$f$. Per tant, el pendent $f'(-1)$ no l'has de calcular: llegeix-lo directament a la gràfica. Quin valor pren la corba a $x = -1$?}

\subsection*{4.2. \normalfont Estudieu si $f(x)$ té punts crítics (punts on $f'(x) = 0$) i, en cas que en tingui, determineu-ne les abscisses i classifiqueu-los segons si són màxims, mínims o punts d'inflexió de $f(x)$. \hspace*{\fill}\textnormal{[1~punt]}}

\pista{Els punts crítics són els punts on $f'(x) = 0$. A la gràfica, són els punts on la corba d'$f'$ talla o toca l'eix d'abscisses.

\medskip
Per classificar-los no cal calcular res: mira el signe d'$f'$ just abans i just després de cada punt.
\begin{itemize}
\item On la gràfica d'$f'$ és per sobre de l'eix d'abscisses, $f'$ és positiva i $f$ creix.
\item On la gràfica d'$f'$ és per sota de l'eix d'abscisses, $f'$ és negativa i $f$ decreix.
\end{itemize}
Llavors, per a cada punt crític:
\begin{itemize}
\item si $f$ passa de créixer a decréixer, hi ha un màxim;
\item si $f$ passa de decréixer a créixer, hi ha un mínim;
\item si $f'$ no canvia de signe, no hi ha ni màxim ni mínim, sinó un punt d'inflexió d'$f$.
\end{itemize}}

\subsection*{4.3. \normalfont Sabent, a més, que $f(0) = -0{,}2$, calculeu l'àrea entre la gràfica de $f'(x)$ i l'eix de les abscisses, des de $x = -1$ fins a $x = 0$. \hspace*{\fill}\textnormal{[0,5~punts]}}

\pista{L'àrea és el valor absolut de la integral d'$f'$:
\[
\text{Àrea} = \left| \int_{-1}^{0} f'(x)\,dx \right|
\]
Mira a la gràfica que $f'$ no canvia de signe entre $x = -1$ i $x = 0$. Per aquest motiu, amb una sola integral ja és suficient.

\medskip
Per calcular-la, aplica la regla de Barrow:
\[
\int_{a}^{b} g(x)\,dx = G(b) - G(a),
\]
on $G$ és una primitiva de $g$. En aquest cas, com que precisament estàs buscant una primitiva d'$f'(x)$, resulta que la mateixa funció $f(x)$ ja et serveix. Per tant, no cal que trobis l'expressió d'$f'$: només necessites $f(0)$, que és a l'enunciat, i $f(-1)$, que surt del punt de l'apartat 4.1. Escriu a la resposta que $f$ és una primitiva d'$f'$, perquè es puntua.}

\end{document}
